% Propietario conservado: /Users/ruben/Documents/ChatGPT/jueces y controles/REVISION_ARGUMENTAL_APERTURAS_CIERRES_20260915/II/manuscript_es/sections/10d_portador_pentadico_rev06.tex
% RESULTADO_RECUPERADO. Integral activo, capítulos 63--65.
% B023_ch63_descenso_icosaedrico_portador_6ba97d5bb029_07_gravedad.tex:196--380.
% B024_ch64_entrelazador_tensorial_reduccion_6ba97d5bb029_07_gravedad.tex:1--122.
% B025_ch65_descomposicion_helicoidal_espin_dos_6ba97d5bb029_07_gravedad.tex:1--104.
% Adecuación editorial: jerarquía, referencias locales y distinción del proyector icosaédrico.
\subsection{Incidencia icosaédrica y acoplamiento gravitatorio pentacomponente}
\label{ii:sec:ii-pentadica}

La estructura discreta del continuo conserva una lectura de incidencia
junto a la lectura aritmética. El portador gravitatorio pentacomponente se
construye mediante esa incidencia y su representación tensorial. La sección
de acoplamiento \(G_a(\vartheta^\circ_{108})\), con
\(a\in\{\mathrm{pre},\mathrm{ret}\}\), es la obtenida en la
sección~\ref{ii:sec:gravity}; a continuación se explicita el espacio sobre el
que actúa, sus cinco componentes y su proyección transversal.

Sea \(X_{\rm TPK}^{\rm exc}\) el subespacio del estado \(\TPK\) enriquecido
que conserva, adem\'as de la salida aritm\'etica, la orientaci\'on, la ruta,
el acarreo y la incidencia excepcional. Para explicitar el espacio geom\'etrico de llegada,
sea \(\varphi_g=(1+\sqrt5)/2\), \(\nu=\sqrt{1+\varphi_g^2}\), y definamos
\begin{equation}
\label{ii:ii:pent:eq:gravity-icosahedron-vertices}
\Omega_{12}^{\rm ico}
=\frac1\nu\bigl(
\{0\}\times\{\pm1\}\times\{\pm\varphi_g\}
\;\cup\;
\{\pm1\}\times\{\pm\varphi_g\}\times\{0\}
\;\cup\;
\{\pm\varphi_g\}\times\{0\}\times\{\pm1\}
\bigr)\subset S^2.
\end{equation}
Sus doce elementos son las clases orientadas de v\'ertice de un icosaedro.
La incidencia no se deja verbal: se fija por
\begin{equation}
\label{ii:ii:pent:eq:gravity-icosahedron-incidence}
\mathcal I_{12}^{\rm ico}
=\{(v,w)\in\Omega_{12}^{\rm ico}\times\Omega_{12}^{\rm ico}:
v\neq w,\ \langle v,w\rangle=1/\sqrt5\}.
\end{equation}
Cada v\'ertice tiene cinco vecinos en esta relaci\'on. El puente de lectura
se tipa como una aplicaci\'on
\begin{equation}
\label{ii:ii:pent:eq:gravity-exceptional-reader-map}
\pi_{\rm ico}:X_{\rm TPK}^{\rm exc}\longrightarrow\Omega_{12}^{\rm ico}.
\end{equation}
La realización de incidencia exige, además de los doce elementos de la
imagen, las siguientes propiedades simultáneas de la aplicación:
\begin{enumerate}[label=(\roman*)]
\item ser sobreyectiva y constante exactamente en las fibras que descartan
la carta aritm\'etica pero conservan la ruta excepcional;
\item entrelazar la acci\'on de los automorfismos orientados del antecedente
con una acci\'on transitiva sobre \(\Omega_{12}^{\rm ico}\);
\item transportar la incidencia TPK declarada a
\(\mathcal I_{12}^{\rm ico}\), no
s\'olo sus cardinales;
\item conservar la involuci\'on de hoja y el retorno non\'adico en el
estabilizador de la fibra correspondiente.
\end{enumerate}

Estas condiciones fijan el orden de la construcción: primero el mapa
\(\pi_{\rm ico}\), después la identificación de la acción inducida y,
finalmente, su representación lineal. La orientación, la ruta y el acarreo
pertenecen al antecedente enriquecido de ese descenso.


% BEGIN OWNER /Users/ruben/Documents/ChatGPT/jueces y controles/REVISION_ARGUMENTAL_APERTURAS_CIERRES_20260915/II/manuscript_es/sections/representacion_dodecafasica_rev08.tex
% FORMALIZACION_REUNIDA de la proyección del calendario y la representación
% de doce posiciones ya desarrolladas en la arquitectura TPK y el enlace Witt.
% Formalización explícita del cotejo de Gram y del transporte de marco.
\subsubsection{Proyección del calendario y representación de sus doce posiciones}
\label{ii:ii:dodeca:calendario}
El calendario de la sección~\ref{ii:ii:tpk-arquitectura} proporciona una
coordenada observable del estado completo. Con sus notaciones, escribimos
\(r_0(\theta)=r(\theta)-1\in\{0,\ldots,8\}\) y
\(b=\beta(\theta)\in\mathbb Z/12\mathbb Z\). La proyección de reloj es
\[
 Q_{\rm clk}(x)=\bigl(\beta(\theta(x)),r_0(\theta(x))\bigr).
\]
Para una ruta admisible de \(n\) actualizaciones, el transporte inducido es
\begin{equation}
 k(r_0,n)=\left\lfloor\frac{r_0+n}{9}\right\rfloor,
 \qquad
 T_n(b,r_0)=\bigl(b+k(r_0,n)\bmod12,\ (r_0+n)\bmod9\bigr).
 \label{ii:ii:dodeca:transporte-reloj}
\end{equation}
Aquí \(n\in\mathbb N_0\) cuenta actualizaciones efectivas; el cociente
entero \(k\) registra los retornos nonádicos atravesados.

\begin{proposition}[Compatibilidad del calendario con la composición de rutas]
\label{ii:ii:dodeca:composicion}
La proyección \(Q_{\rm clk}\), junto con la longitud de cada ruta, define
un funtor de la categoría de rutas admisibles TPK a la categoría de acción
del calendario~\eqref{ii:ii:dodeca:transporte-reloj}. Se cumplen
\(T_{m+n}=T_mT_n\), \(T_0=I\),
\(T_9(b,r_0)=(b+1,r_0)\) y \(T_{108}(b,r_0)=(b,r_0)\).
\end{proposition}
\begin{proof}
Poniendo \(r'_0=(r_0+n)\bmod9\), la división euclidiana da
\[
 k(r_0,n+m)=k(r_0,n)+k(r'_0,m).
\]
Esta identidad conserva simultáneamente residuo y cociente al concatenar
las rutas. La actualización de \(\theta\) establecida en la
sección~\ref{ii:ii:tpk-arquitectura} induce exactamente
\eqref{ii:ii:dodeca:transporte-reloj}; las identidades restantes se obtienen
por sustitución. La categoría de llegada tiene los pares \((b,r_0)\) como
objetos y las longitudes admisibles como flechas, con suma como composición.
\end{proof}

El índice \(p=b+1\in\{1,\ldots,12\}\), tomando el representante
\(0\le b<12\), publica la posición dodecafásica. Una vuelta de
\(108\) actualizaciones recorre sus doce posiciones. La periodicidad
demostrada corresponde a esta proyección: ruta, acarreo acumulado y
memoria permanecen en el estado fuente y continúan su actualización.

\paragraph{Realización de incidencia a partir del proyector angular.}
Los representantes \(r_pC_5\) y la representación \(\rho\) de la
sección~\ref{ii:sec:ii-enlace-angular-witt} fijan la carta de las doce
posiciones. El proyector ya definido por
\eqref{ii:eq:ii-witt-projector-character} produce
\begin{equation}
 \begin{gathered}
 \mathcal H_3=P_3^{\mathrm{ico}}\mathbb R^{12},\qquad v_p=2P_3^{\mathrm{ico}}e_p,
 \\[-2pt]
 (P_3^{\mathrm{ico}})_{pq}=\frac1{20}
 \sum_{\substack{g\in A_5\\gr_qC_5=r_pC_5}}\chi_3(g^{-1}).
 \end{gathered}
 \label{ii:ii:dodeca:vertices-proyector}
\end{equation}
La última suma contiene cinco términos por entrada y utiliza los
representantes y las dos clases de ciclos de orden cinco allí fijados.

\begin{proposition}[Realización icosaédrica de las posiciones dodecafásicas]
\label{ii:ii:dodeca:realizacion}
Los vectores \(v_p\) son doce vectores unitarios distintos de
\(\mathcal H_3\). Su incidencia es
\(p\sim q\Longleftrightarrow\langle v_p,v_q\rangle=1/\sqrt5\).
La aplicación \(p\mapsto v_p\) es \(A_5\)-equivariante y su imagen
es isométrica a \(\Omega_{12}^{\rm ico}\).
\end{proposition}
\begin{proof}
La transitividad y \(\operatorname{Tr}P_3^{\mathrm{ico}}=3\) dan
\((P_3^{\mathrm{ico}})_{pp}=1/4\); por idempotencia,
\(\langle v_p,v_q\rangle=4(P_3^{\mathrm{ico}})_{pq}\) y \(\|v_p\|=1\).
La evaluación de los cinco términos de
\eqref{ii:ii:dodeca:vertices-proyector} da los pares antipodales
\[
 (1,4),\ (2,3),\ (5,8),\ (6,7),\ (9,12),\ (10,11).
\]
Para los representantes \((1,2,5,6,9,10)\), su matriz de Gram es
\begin{equation}
 I_6+\frac D{\sqrt5},\qquad
 D=\begin{pmatrix}
 0&-1&-1&1&-1&1\\
 -1&0&1&1&-1&-1\\
 -1&1&0&1&1&1\\
 1&1&1&0&-1&1\\
 -1&-1&1&-1&0&1\\
 1&-1&1&1&1&0
 \end{pmatrix},\qquad D^2=5I_6.
 \label{ii:ii:dodeca:gram}
\end{equation}
Los restantes productos escalares se obtienen mediante antipodía.
Así, aparte de un vértice y su opuesto, los productos son
\(\pm1/\sqrt5\), y cada vértice tiene cinco vecinos. La tabla siguiente
explicita el cotejo con la carta firmada de seis posiciones
\(\mathcal U=(\infty,0,1,2,3,4)\):
\begin{equation}
\begin{array}{c|cc@{\qquad}c|cc}
p&u(p)&\sigma(p)&p&u(p)&\sigma(p)\\ \hline
1&\infty&+1&7&3&-1\\
2&0&-1&8&1&+1\\
3&0&+1&9&2&-1\\
4&\infty&-1&10&4&+1\\
5&1&-1&11&4&-1\\
6&3&+1&12&2&+1
\end{array}
\label{ii:ii:dodeca:carta-firmada}
\end{equation}
En efecto, para \(C=\operatorname{bal}(A_W)\), con \(A_W\) dado en
\eqref{ii:exc:aw}, las entradas de \(D\) son
\(D_{ab}=\sigma(p_a)\sigma(p_b)C_{u(p_a),u(p_b)}\).
Por consiguiente, la tabla y~\eqref{ii:ii:dodeca:gram} identifican
isométricamente \(v_p\) con
\(\sigma(p)\sqrt2 E_Ce_{u(p)}\), donde
\(E_C=(I_6+C/\sqrt5)/2\). Ambos sistemas generan espacios de dimensión
tres y tienen idéntica matriz de Gram, lo que prueba que la aplicación
lineal inducida está bien definida y es isométrica. La realización
explícita de~\eqref{ii:ii:ico:proyector-seis}, demostrada a continuación,
la identifica con~\(\Omega_{12}^{\rm ico}\) y su incidencia.
Finalmente, \(P_3^{\mathrm{ico}}\rho(g)=\rho(g)P_3^{\mathrm{ico}}\) implica
\(v_{g\cdot p}=\rho(g)v_p\). La biyección de las doce posiciones da
\(\operatorname{Stab}_{A_5}(v_p)=r_pC_5r_p^{-1}\).
\end{proof}

Se obtiene así la composición representacional
\begin{equation}
 x\longmapsto Q_{\rm clk}(x)\longmapsto
 p=\beta(\theta(x))+1\longmapsto v_p.
 \label{ii:ii:dodeca:composicion-representacional}
\end{equation}
La expresión firmada de esta composición es
\(r_\pm^{\rm rep}(x)=(u(p),\sigma(p))\), determinada por la
tabla~\eqref{ii:ii:dodeca:carta-firmada}. La tabla es una elección explícita
de alineamiento entre cartas; las acciones de \(A_5\) proporcionan sus
alineamientos equivalentes. El registro completo se conserva en el
dominio de~\eqref{ii:ii:dodeca:composicion-representacional}.

\paragraph{Enlace con el portador tensorial de cinco componentes.}
En \(\mathcal H_3\), formamos
\(Q_p=v_pv_p^*-I_{\mathcal H_3}/3\). Sea \(J=\mathbf1\mathbf1^*\)
y sea \(\mathsf A\) la matriz de la permutación antipodal anterior.
El producto de Frobenius y~\eqref{ii:ii:dodeca:gram} dan
\begin{equation}
 \begin{aligned}
 \operatorname{Gram}(Q_p)
 &=(4P_3^{\mathrm{ico}})\circ(4P_3^{\mathrm{ico}})-\frac13J
 =\frac45(I+\mathsf A)-\frac2{15}J,
 \\
 P^{\rm ico}_5&=\frac58\operatorname{Gram}(Q_p)
 =\frac12(I+\mathsf A)-\frac1{12}J.
 \end{aligned}
 \label{ii:ii:dodeca:puente-tensorial}
\end{equation}
El símbolo \(\circ\) denota el producto entrada a entrada. La última
expresión es el proyector ortogonal sobre las coordenadas pares bajo
antipodía y ortogonales a \(\mathbf1\), de dimensión cinco. Su carácter
es \((5,1,-1,0,0)\): sobre los seis pares antipodales, los números de
puntos fijos son \((6,2,0,1,1)\), y se retira la componente constante.
Por ello coincide con el proyector central de
\eqref{ii:ii:pent:eq:gravity-p5}. Esta identidad une la construcción
tridimensional con \(\Gamma_0\) y \(\Gamma_5\), definidos en
\eqref{ii:ii:pent:eq:gravity-Gamma0} y~\eqref{ii:ii:pent:eq:gravity-gamma5},
mediante la operación cuadrática explícita \(v_p\mapsto Q_p\).

\subsubsection{Transporte del marco representacional}
\label{ii:ii:dodeca:marco}
Sea \(s=(1\,2\,\cdots\,12)\) y sea \(S_ce_p=e_{s(p)}\).
Un retorno nonádico del calendario lleva \(p\) a \(s(p)\). El transporte
conjunto de marco, representación y proyector se expresa por
\begin{equation}
 \begin{gathered}
 P_{3,k}^{\mathrm{ico}}=S_c^kP_3^{\mathrm{ico}}S_c^{-k},\qquad
 \rho^{(k)}(g)=S_c^k\rho(g)S_c^{-k},\\
 v^{(k)}_{s^k(p)}=2P_{3,k}^{\mathrm{ico}}e_{s^k(p)}=S_c^kv_p.
 \end{gathered}
 \label{ii:ii:dodeca:marco-movil}
\end{equation}
La última igualdad se prueba por sustitución; la ortogonalidad de
\(S_c\) conserva todos los productos escalares y transporta la
incidencia. Para una ruta de longitud \(n\), se toma
\(k=k(r_0,n)\). La identidad de acarreo de la
proposición~\ref{ii:ii:dodeca:composicion} prueba la compatibilidad de
estos transportes con la concatenación. Aunque \(S_c^{12}=I\),
las doce vueltas actualizan la memoria de la ruta fuente.

La igualdad~\eqref{ii:ii:dodeca:marco-movil} establece covariancia entre
marcos. La acción activa sobre una incidencia fija tiene un contenido
adicional: requiere identificar el transporte del estado completo y su
acción en esa carta. En particular, \(s\) tiene orden doce, mientras
que los órdenes de los elementos de \(A_5\) son \(1,2,3,5\).
Asimismo, el estabilizador de un vértice tiene orden cinco, y el de su
par antipodal tiene orden diez; la involución de la carta firmada
intercambia los dos vértices. Estas distinciones conservan el significado
de orientación, hoja y retorno en las propiedades (i)--(iv): la
realización representacional demostrada y el descenso de la ruta
excepcional completa se relacionan mediante esas propiedades, con sus
dominios y fibras declarados.

% END OWNER


% BEGIN OWNER /Users/ruben/Documents/ChatGPT/jueces y controles/REVISION_ARGUMENTAL_APERTURAS_CIERRES_20260915/II/manuscript_es/sections/incidencia_icosaedrica_rev07.tex
% FORMALIZACION_NUEVA de la incidencia firmada de la carta Paley preexistente.
% La construcción comienza en exc:aw; no identifica por decreto sus banderas
% con el cociente de todas las historias TPK.
\paragraph{Construcción del portador desde la incidencia de seis posiciones.}
\label{ii:ii:ico:construccion-firmada}
La matriz \(A_W\) de \eqref{ii:exc:aw}, que interviene previamente en el
transporte angular de la sección~\ref{ii:sec:ii-enlace-angular-witt}, proporciona
un antecedente explícito de la realización icosaédrica. Sea
\(C=\operatorname{bal}(A_W)\), donde
\(\operatorname{bal}(0)=0\), \(\operatorname{bal}(1)=1\) y
\(\operatorname{bal}(2)=-1\). El orden de sus coordenadas es
\(\mathcal U=(\infty,0,1,2,3,4)\). La carta entera satisface
\(C=C^{\mathsf T}\), \(C^2=5I_6\) y \(\operatorname{Tr}C=0\).
Definimos su recubrimiento firmado por
\begin{equation}
 \mathcal D_C=\mathcal U\times\{+1,-1\},\qquad
 (u,\sigma)\sim_C(v,\tau)
 \ \Longleftrightarrow\ u\ne v\ \text{ y }\ \sigma\tau=C_{uv}.
 \label{ii:ii:ico:incidencia-firmada}
\end{equation}
La orientación duplica cada coordenada; la matriz decide las incidencias
entre las dos fibras correspondientes. La involución de este recubrimiento
es \(\jmath_C(u,\sigma)=(u,-\sigma)\).

\begin{proposition}[Realización isométrica de la incidencia firmada]
\label{ii:ii:ico:realizacion-firmada}
El operador y los vectores
\begin{equation}
 E_C=\frac12\left(I_6+\frac C{\sqrt5}\right),\qquad
 \mathcal H_C=\operatorname{im}E_C,\qquad
 v_{u,\sigma}=\sigma\sqrt2 E_Ce_u
 \label{ii:ii:ico:proyector-seis}
\end{equation}
realizan \(\mathcal D_C\) como doce vectores unitarios distintos en un
espacio euclidiano de dimensión tres. Existe una isometría explícita
\(T_C:\mathcal H_C\to\mathbb R^3\) para la que
\(\iota_C(u,\sigma)=T_Cv_{u,\sigma}\) es una biyección sobre
\(\Omega_{12}^{\rm ico}\), conserva la incidencia y satisface
\(\iota_C\jmath_C=-\iota_C\).
\end{proposition}
\begin{proof}
De \(C^2=5I_6\) se obtiene \(E_C^2=E_C=E_C^*\); su traza es tres.
Las identidades métricas son
\[
 \langle v_{u,\sigma},v_{v,\tau}\rangle
 =\begin{cases}\sigma\tau,&u=v,\\
 \sigma\tau C_{uv}/\sqrt5,&u\ne v.
 \end{cases}
\]
Así, los doce vectores son distintos y su incidencia de producto escalar
\(1/\sqrt5\) es exactamente~\eqref{ii:ii:ico:incidencia-firmada}.
Para explicitar la isometría, sean las columnas de \(B\), en el orden
\(\mathcal U\),
\[
 B=\frac1\nu
 \begin{pmatrix}
 0&-1&-\varphi_g&0&\varphi_g&1\\
 -1&-\varphi_g&0&1&0&-\varphi_g\\
 -\varphi_g&0&-1&-\varphi_g&-1&0
 \end{pmatrix}.
\]
La sustitución de \(\varphi_g^2=\varphi_g+1\) da
\(B^*B=I_6+C/\sqrt5=2E_C\). Por tanto
\(T_C=B/\sqrt2\big|_{\mathcal H_C}\) es una isometría, y
\(T_Cv_{u,\sigma}=\sigma Be_u\). Las seis columnas de \(B\) y sus
opuestas son precisamente los doce vértices de
\eqref{ii:ii:pent:eq:gravity-icosahedron-vertices}. La última identidad
demuestra también la compatibilidad de las involuciones.
\end{proof}

La realización dodecafásica precedente construye los doce vectores
directamente desde el proyector \(P_3^{\mathrm{ico}}\) y las posiciones del calendario.
El recubrimiento firmado proporciona otra expresión coordenada de ese
portador. Para el lector \(\pi_{\rm ico}\) del teorema de descenso,
la relación entre ambas expresiones se escribe
\begin{equation}
 X_{\rm TPK}^{\rm exc}\xrightarrow{\ r_\pm\ }
 \mathcal D_C\xrightarrow[\simeq]{\ \iota_C\ }
 \Omega_{12}^{\rm ico},\qquad
 \pi_{\rm ico}(x)=\sigma(x)Be_{u(x)},\quad
 r_\pm(x)=(u(x),\sigma(x)).
 \label{ii:ii:ico:factorizacion-lector}
\end{equation}
Como \(\iota_C\) es una biyección, esta factorización determina
\(r_\pm=\iota_C^{-1}\pi_{\rm ico}\). Por tanto, \(r_\pm\) expresa
en la carta firmada el lector que se esté utilizando: su introducción
no añade una hipótesis independiente para construir el portador
dodecafásico. La tabla de correspondencias precedente da esa expresión
para sus doce posiciones. Las propiedades (i)--(iv) conservan su
alcance propio cuando se identifica la realización con el cociente
de la ruta excepcional completa. En particular, el transporte conjunto
de marco, proyector e incidencia demuestra la covariancia entre cartas;
la conservación de una incidencia fija por una transformación activa
es una identidad distinta. La orientación y la memoria permanecen
en el estado fuente durante ambas lecturas.
Una permutación firmada que transporta \(C\) transporta conjuntamente
\(E_C\), sus vectores y su incidencia; esta covariancia permite cotejar
las cartas conservando la orientación de cada fibra.

% END OWNER


\begin{theorem}[Descenso de incidencia al cociente icosa\'edrico]
\label{ii:ii:pent:thm:gravity-exceptional-descent}
Sup\'ongase que \(\pi_{\rm ico}\) satisface (i)--(iv) y que la acci\'on
orientada inducida en \(\Omega_{12}^{\rm ico}\) tiene grupo
\(G_{\rm ico}\simeq A_5\). Entonces, para toda clase
\(v_0\in\Omega_{12}^{\rm ico}\),
\begin{equation}
\label{ii:ii:pent:eq:gravity-icosahedral-quotient}
H_{v_0}=\operatorname{Stab}_{G_{\rm ico}}(v_0)\simeq C_5,
\qquad
\Omega_{12}^{\rm ico}\simeq G_{\rm ico}/H_{v_0}\simeq A_5/C_5.
\end{equation}
La aplicaci\'on
\begin{equation}
\label{ii:ii:pent:eq:gravity-coset-map}
\kappa_{v_0}:G_{\rm ico}/H_{v_0}\longrightarrow\Omega_{12}^{\rm ico},
\qquad gH_{v_0}\longmapsto g\cdot v_0,
\end{equation}
es una biyecci\'on equivariante y preserva la incidencia si en el cociente
se declara
\begin{equation}
\label{ii:ii:pent:eq:gravity-coset-incidence}
gH_{v_0}\sim g'H_{v_0}
\quad\Longleftrightarrow\quad
(g\cdot v_0,g'\cdot v_0)\in\mathcal I_{12}^{\rm ico}.
\end{equation}
Por tanto, el compuesto \(\kappa_{v_0}^{-1}\pi_{\rm ico}\) es el mapa
tipado del antecedente TPK al cociente \(A_5/C_5\).
\end{theorem}

\begin{proof}
La lista de \eqref{ii:ii:pent:eq:gravity-icosahedron-vertices} tiene doce elementos.
Por transitividad, \(|G_{\rm ico}\cdot v_0|=12\), y el teorema
\`orbita--estabilizador da \(|H_{v_0}|=60/12=5\). Todo grupo de orden
primo cinco es c\'iclico; por tanto, \(H_{v_0}\simeq C_5\). La f\'ormula de
\(\kappa_{v_0}\) no depende del representante: si
\(gH_{v_0}=g'H_{v_0}\), entonces \(g^{-1}g'\in H_{v_0}\) y
\(g'v_0=gv_0\). La sobreyectividad procede de la transitividad, y la
igualdad de cardinales da inyectividad. La equivarianza es inmediata:
\(\kappa_{v_0}(agH_{v_0})=a\kappa_{v_0}(gH_{v_0})\). Finalmente,
\eqref{ii:ii:pent:eq:gravity-coset-incidence} es precisamente el transporte de
estructura por esa biyecci\'on; la hip\'otesis (iii) hace conmutar este
transporte con \(\pi_{\rm ico}\).
\end{proof}

El teorema establece el descenso una vez verificadas las cuatro
propiedades de \(\pi_{\rm ico}\). Su dominio es el estado enriquecido
de la lectura excepcional, que conserva la incidencia transportada por
APP--TRIT--TPK. En esta construcción, el antecedente explícito es la
aplicación~\eqref{ii:ii:pent:eq:gravity-exceptional-reader-map} con sus cuatro
propiedades; el tramo desde \(A_5/C_5\) hasta la reducción transversal sin
traza se demuestra a continuación con ese antecedente y la acción inducida
de \(A_5\).

La cadena de procedencia se expresa mediante todos sus morfismos:
\begin{equation}
\label{ii:ii:pent:eq:gravity-exceptional-provenance}
\APP\longrightarrow\mathrm{TRIT}\longrightarrow\TPK
\longrightarrow X_{\rm TPK}^{\rm exc}
\xrightarrow{\ \pi_{\rm ico}\ }\Omega_{12}^{\rm ico}
\xleftarrow[\simeq]{\ \kappa_{v_0}\ }A_5/C_5.
\end{equation}
La composición~\eqref{ii:ii:pent:eq:gravity-exceptional-provenance} conserva
como antecedente la existencia de \(\pi_{\rm ico}\) con las propiedades
(i)--(iv). Sobre ese dominio, el descenso al cociente es el
teorema~\ref{ii:ii:pent:thm:gravity-exceptional-descent}; de él se construyen
el proyector \(P^{\mathrm{ico}}_5\), el entrelazador \(\Gamma_5\) y la
proyección \(\Lambda(n)\). La realización del portador como campo
gravitatorio utiliza, además, la dinámica y las restricciones físicas que
se especifican al final de este desarrollo.

Sea \(\mathcal G=A_5\), sea \(H\simeq C_5\) el estabilizador de un v\'ertice
orientado del icosaedro y consid\'erese
\begin{equation}
\label{ii:ii:pent:eq:gravity-v12}
V_{12}=\R[\mathcal G/H],
\qquad |\mathcal G/H|=\frac{60}{5}=12.
\end{equation}

% BEGIN OWNER /Users/ruben/Documents/ChatGPT/jueces y controles/REVISION_ARGUMENTAL_APERTURAS_CIERRES_20260915/II/manuscript_es/sections/caracteres_pentadicos_rev07.tex
% FORMALIZACION_NUEVA: prueba explícita de la descomposición y la
% irreducibilidad ya enunciadas en el desarrollo pentacomponente de REV06.
% CERTIFICADO_NUEVO: verificar_caracteres_pentadicos_rev07.py.
% Inserción prevista tras la definición de V12, antes de su descomposición.
\paragraph{Caracteres de la incidencia icosaédrica.}
El cálculo siguiente se efectúa sobre la acción icosaédrica obtenida en el
teorema~\ref{ii:ii:pent:thm:gravity-exceptional-descent}, con sus hipótesis
de incidencia y orientación. Determina la descomposición de la
representación de vértices y la irreducibilidad del espacio tensorial
mediante sus trazas, conservando el antecedente
\(\pi_{\mathrm{ico}}\) declarado en ese teorema.

\begin{proposition}[Descomposición de vértices y carácter tensorial]
\label{ii:prop:ii-pent-caracteres}
Sean \(\tau=(1+\sqrt5)/2\) y
\(\bar\tau=(1-\sqrt5)/2\). En el orden de clases
\(1,2A,3A,5A,5B\), los caracteres de la representación rotacional
\(V_3\), de su conjugada algebraica \(V_{3'}\), de la representación
de los doce vértices \(V_{12}\) y de
\(\operatorname{STF}_3=\operatorname{Sym}^2_0(V_3)\) son
\begin{equation}
\label{ii:eq:ii-pent-tabla-caracteres}
\begin{array}{c|rrrrr}
\text{clase} &1&2A&3A&5A&5B\\ \hline
\text{cardinal}&1&15&20&12&12\\
\text{clase de }g^2&1&1&3A&5B&5A\\ \hline
\chi_3&3&-1&0&\tau&\bar\tau\\
\chi_{3'}&3&-1&0&\bar\tau&\tau\\
\chi_{12}&12&0&0&2&2\\
\chi_{\mathrm{STF}}&5&1&-1&0&0
\end{array}
\end{equation}
Las representaciones \(V_3,V_{3'}\) y
\(\operatorname{STF}_3\) son absolutamente irreducibles y distintas.
Además,
\begin{equation}
\label{ii:eq:ii-pent-descomposicion-demostrada}
 V_{12}\simeq\mathbf1\oplus V_3\oplus V_{3'}
                  \oplus\operatorname{STF}_3,
 \qquad
 \langle\chi_{12},\chi_{\mathrm{STF}}\rangle=1.
\end{equation}
En particular, el sumando de dimensión cinco aparece una sola vez.
\end{proposition}

\begin{proof}
En \(A_5\), las dobles transposiciones son \(5\cdot3=15\) y
los ciclos de longitud tres son \(\binom53\cdot2=20\).
Los \(4!=24\) ciclos de longitud cinco se separan en dos clases de
doce: su centralizador es el subgrupo cíclico de orden cinco y está
contenido en \(A_5\). Si se identifican las posiciones de un ciclo
con \(\mathbb Z/5\mathbb Z\), la permutación que induce
\(r\mapsto2r\) es un ciclo de longitud cuatro sobre los elementos
no nulos, de signo negativo. Por ello el cuadrado intercambia ambas
clases. La inversión \(r\mapsto-r\) es producto de dos
transposiciones y conserva cada clase. Esto determina las dos primeras
filas de~\eqref{ii:eq:ii-pent-tabla-caracteres}.

La representación \(V_3\) es la acción por rotaciones sobre las
coordenadas de los vértices. Una rotación de ángulo \(\theta\)
tiene traza \(1+2\cos\theta\). Las clases de órdenes dos y tres
dan, respectivamente, \(-1\) y \(0\). Se denomina \(5A\) a la
clase de ángulos \(\pm2\pi/5\) y \(5B\) a la de
\(\pm4\pi/5\); sus trazas son \(\tau\) y \(\bar\tau\).
Estas expresiones se obtienen de
\(u^2+u-1=0\), para \(u=2\cos(2\pi/5)\), y de la fórmula
del ángulo doble. La ecuación de \(u\) resulta al dividir
\(1+z+z^2+z^3+z^4=0\) por \(z^2\), con
\(z=\exp(2\pi i/5)\). Se ha calculado así el carácter de una
representación ya definida geométricamente.

Para construir la otra representación, se usan los vértices sin la
normalización común \(\nu^{-1}\). Sus coordenadas pertenecen a
\(K=\mathbb Q(\sqrt5)\). Si \(B\) es la matriz de tres vértices
linealmente independientes y \(B_g\) contiene sus imágenes, la
matriz rotacional es \(R_g=B_gB^{-1}\in M_3(K)\).
La conjugación \(\sigma:\sqrt5\mapsto-\sqrt5\) define entonces
\(R'_g=\sigma(R_g)\), y satisface
\[
 R'_{gh}=R'_gR'_h,\qquad
 (R'_g)^TR'_g=I,\qquad \det R'_g=1.
\]
Por tanto, \(V_{3'}\) es una representación real efectivamente
construida. Su carácter es \(\sigma(\chi_3)\), que intercambia
\(\tau\) y \(\bar\tau\).

El carácter de permutación cuenta vértices fijos. La identidad fija
doce; las rotaciones no triviales de orden cinco fijan los dos extremos
de su eje. Las de orden dos o tres no fijan vértices, pues el
estabilizador de un vértice es \(C_5\). Se obtiene así
\(\chi_{12}=(12,0,0,2,2)\).

Para el espacio tensorial, si \(\lambda_1,\lambda_2,\lambda_3\)
son los autovalores de \(R_g\), la suma
\(\sum_{i\leq j}\lambda_i\lambda_j\) es
\(\bigl((\sum_i\lambda_i)^2+\sum_i\lambda_i^2\bigr)/2\).
La forma métrica genera una recta invariante en
\(\operatorname{Sym}^2(V_3)\), cuyo complemento sin traza tiene
por carácter
\begin{equation}
\label{ii:eq:ii-pent-caracter-stf}
 \chi_{\mathrm{STF}}(g)
 =\frac{\chi_3(g)^2+\chi_3(g^2)}2-1.
\end{equation}
Al usar el mapa de cuadrados de la tabla y
\(\tau+\bar\tau=1\), \(\tau\bar\tau=-1\), resulta
\(\chi_{\mathrm{STF}}=(5,1,-1,0,0)\).

El producto de caracteres reales es
\(\langle\chi,\psi\rangle
=\frac1{60}\sum_C|C|\chi(C)\psi(C)\).
En particular,
\begin{align*}
 \langle\chi_3,\chi_3\rangle
 &=\langle\chi_{3'},\chi_{3'}\rangle
 =\frac{9+15+12(\tau^2+\bar\tau^2)}{60}=1,\\
 \langle\chi_3,\chi_{3'}\rangle
 &=\frac{9+15+24\tau\bar\tau}{60}=0,\\
 \langle\chi_{\mathrm{STF}},\chi_{\mathrm{STF}}\rangle
 &=\frac{25+15+20}{60}=1.
\end{align*}
La norma uno prueba irreducibilidad de las complejificaciones; por
consiguiente, también irreducibilidad absoluta de las representaciones
reales construidas. Sus dimensiones y la ortogonalidad distinguen los
tres sumandos. Finalmente, la igualdad de funciones de clase
\[
 \chi_{12}=1+\chi_3+\chi_{3'}+\chi_{\mathrm{STF}}
\]
se verifica en las cinco columnas. La semisimplicidad de las
representaciones de un grupo finito en característica cero da la
descomposición~\eqref{ii:eq:ii-pent-descomposicion-demostrada}. La
multiplicidad del sumando tensorial puede comprobarse directamente:
\[
 \langle\chi_{12},\chi_{\mathrm{STF}}\rangle
 =\frac{12\cdot5}{60}=1.
\]
\end{proof}

Así, el idempotente central asociado a
\(\chi_{\mathrm{STF}}=\chi_5\) tiene rango
\(5\cdot1=5\) en \(V_{12}\). La irreducibilidad que interviene
en la normalización del entrelazador tensorial y la multiplicidad que
interviene en el acoplamiento gravitatorio quedan acreditadas por
este mismo cálculo.

% END OWNER


La representaci\'on de permutaci\'on \(\rho_{12}\) se descompone como
\begin{equation}
\label{ii:ii:pent:eq:gravity-v12-decomposition}
V_{12}\simeq\mathbf1\oplus V_3\oplus V_{3'}\oplus V_5.
\end{equation}
El idempotente central del sumando de car\'acter
\(\chi_5=(5,1,-1,0,0)\) es
\begin{equation}
\label{ii:ii:pent:eq:gravity-p5}
\boxed{
P^{\mathrm{ico}}_5=\frac5{60}\sum_{g\in A_5}\chi_5(g^{-1})\rho_{12}(g)
=\frac1{12}\left(
5I+\sum_{g\in2A}\rho_{12}(g)
-\sum_{g\in3A}\rho_{12}(g)
\right).}
\end{equation}

\begin{theorem}[Proyector y acoplamiento gravitatorio de rango cinco]
\label{ii:ii:pent:thm:gravity-p5-coupling}
El operador \(P^{\mathrm{ico}}_5\) satisface
\begin{equation}
\label{ii:ii:pent:eq:gravity-p5-properties}
(P^{\mathrm{ico}}_5)^2=P^{\mathrm{ico}}_5=(P^{\mathrm{ico}}_5)^*,
\qquad \Tr P^{\mathrm{ico}}_5=\operatorname{rank} P^{\mathrm{ico}}_5=5.
\end{equation}
Para
\begin{equation}
\label{ii:ii:pent:eq:gravity-G5}
\mathbb G_{5,a}:=G_a(\vartheta^\circ_{108})P^{\mathrm{ico}}_5
\end{equation}
se tiene
\begin{equation}
\label{ii:ii:pent:eq:gravity-G5-spectrum}
\Spec(\mathbb G_{5,a})
=\{G_a(\vartheta^\circ_{108})\text{ con multiplicidad }5,
\;0\text{ con multiplicidad }7\}.
\end{equation}
\end{theorem}

\begin{proof}
La ortogonalidad de caracteres convierte \eqref{ii:ii:pent:eq:gravity-p5} en el
idempotente central que act\'ua como la identidad sobre la \`unica copia de
\(V_5\) y como cero en los otros tres sumandos de
\eqref{ii:ii:pent:eq:gravity-v12-decomposition}. La realidad del car\'acter y la
ortogonalidad de \(\rho_{12}\) dan la autoadjunci\'on. Su traza y su rango
son cinco. Multiplicar un proyector de rango cinco por el escalar positivo
\(G_a\) produce inmediatamente el espectro indicado.
\end{proof}

La sección gravitatoria \(G_a\) es común a las cinco componentes del
sumando tensorial seleccionado. Su multiplicidad espectral expresa el
rango del portador; los otros siete modos del módulo de permutación
pertenecen al núcleo de \(\mathbb G_{5,a}\). El símbolo
\(P^{\mathrm{ico}}_5\) identifica aquí este proyector central de rango
cinco sobre \(V_{12}\); el promedio de las cinco posiciones de una
pentafibra aritmética tiene otro dominio y rango uno.

\subsection{Entrelazador con el espacio de tensores sim\'etricos sin traza}

Sea \(\mathcal V=\Omega_{12}^{\rm ico}\subset S^2\) el conjunto de los
doce vértices unitarios ya construido. Para cada \(v\in\mathcal V\), def\'inase
\begin{equation}
\label{ii:ii:pent:eq:gravity-Qv}
Q_v=vv^T-\frac13I_3\in
\operatorname{STF}_3:=\operatorname{Sym}^2_0(\R^3).
\end{equation}
El mapa de marco
\begin{equation}
\label{ii:ii:pent:eq:gravity-Gamma0}
\Gamma_0\!\left(\sum_{v\in\mathcal V}x_ve_v\right)
=\sum_{v\in\mathcal V}x_vQ_v
\end{equation}
es \(A_5\)-equivariante y satisface
\begin{equation}
\label{ii:ii:pent:eq:gravity-tight-frame}
\Gamma_0\Gamma_0^*=\frac85I_{\operatorname{STF}_3},
\qquad
\Gamma_0=\Gamma_0P^{\mathrm{ico}}_5.
\end{equation}

\begin{theorem}[Entrelazador tensorial normalizado]
\label{ii:ii:pent:thm:gravity-gamma5}
La aplicaci\'on
\begin{equation}
\label{ii:ii:pent:eq:gravity-gamma5}
\boxed{
\Gamma_5:=\sqrt{\frac58}\,\Gamma_0|_{V_5}:
V_5\xrightarrow{\ \simeq\ }\operatorname{STF}_3}
\end{equation}
es una isometr\'ia \(A_5\)-equivariante. Su adjunto es
\begin{equation}
\label{ii:ii:pent:eq:gravity-gamma5-adjoint}
\Gamma_5^*Q=\sqrt{\frac58}
\sum_{v\in\mathcal V}(v^TQv)e_v.
\end{equation}
\end{theorem}

\begin{proof}
La equivarianza sigue de \(Q_{gv}=gQ_vg^T\). El operador
\(\Gamma_0\Gamma_0^*\) conmuta con la representaci\'on irreducible de
\(A_5\) sobre \(\operatorname{STF}_3\), por lo que es escalar. Su traza
es
\[
\sum_{v\in\mathcal V}\|Q_v\|_F^2
=12\cdot\frac23=8;
\]
al dividir por la dimensi\'on cinco se obtiene \(8/5\). De aqu\'i,
\(\Gamma_5\Gamma_5^*=I\). Como dominio y codominio tienen dimensi\'on
cinco, tambi\'en \(\Gamma_5^*\Gamma_5=I\). La f\'ormula del adjunto
procede de \(\langle Q,Q_v\rangle_F=v^TQv\).
\end{proof}

La normalización mediante el marco icosaédrico construye la isometría que
transporta el sumando finito al espacio tensorial real de espín dos.

\subsection{Proyecci\'on transversal sin traza}

Para \(n\in S^2\), sea \(\Pi(n)=I-nn^T\). Definimos
\begin{equation}
\label{ii:ii:pent:eq:gravity-lambda-tt}
\boxed{
\Lambda(n)Q
=\Pi Q\Pi-\frac12\Tr(\Pi Q\Pi)\Pi.}
\end{equation}

\begin{theorem}[Reducci\'on de cinco componentes a dos]
\label{ii:ii:pent:thm:gravity-tt}
La restricci\'on de \(\Lambda(n)\) a \(\operatorname{STF}_3\) es el
proyector ortogonal sobre
\begin{equation}
\label{ii:ii:pent:eq:gravity-htt}
\mathcal H_{\rm TT}(n)
=\{Q\in\operatorname{STF}_3:Qn=0\}.
\end{equation}
En particular,
\begin{equation}
\label{ii:ii:pent:eq:gravity-lambda-properties}
\Lambda(n)^2=\Lambda(n)=\Lambda(n)^*,
\qquad \operatorname{rank}\Lambda(n)=2.
\end{equation}
Para \(n=e_3\), la imagen est\'a generada por
\begin{equation}
\label{ii:ii:pent:eq:gravity-plus-cross}
E_+=\frac1{\sqrt2}\operatorname{diag}(1,-1,0),
\qquad
E_\times=\frac1{\sqrt2}
\begin{pmatrix}0&1&0\\1&0&0\\0&0&0\end{pmatrix}.
\end{equation}
\end{theorem}

\begin{proof}
Como \(\Pi n=0\), la salida es transversal. Su traza vale
\[
\Tr(\Pi Q\Pi)-\frac12\Tr(\Pi)\Tr(\Pi Q\Pi)=0,
\]
pues \(\Tr\Pi=2\). Si \(Qn=0\) y \(\Tr Q=0\), entonces
\(\Pi Q\Pi=Q\) y \(\Tr(\Pi Q\Pi)=0\), por lo que
\(\Lambda(n)Q=Q\). La f\'ormula es autoadjunta para el producto de
Frobenius. En el marco \(n=e_3\) su matriz en una base STF adaptada es
\(\operatorname{diag}(1,1,0,0,0)\), lo que prueba el rango.
\end{proof}

La cadena completa queda, por tanto,
\begin{equation}
\label{ii:ii:pent:eq:gravity-12-5-5-2}
\boxed{
V_{12}\xrightarrow{P^{\mathrm{ico}}_5}V_5
\xrightarrow{\Gamma_5}\operatorname{STF}_3
\xrightarrow{\Lambda(n)}\mathcal H_{\rm TT}(n),
\qquad 12\longrightarrow5\longrightarrow5\longrightarrow2.}
\end{equation}
El acoplamiento reducido
\begin{equation}
\label{ii:ii:pent:eq:gravity-effective-map}
\mathbb G_{{\rm TT},a}(n)
=G_a(\vartheta^\circ_{108})\Lambda(n)\Gamma_5P^{\mathrm{ico}}_5
\end{equation}
tiene rango dos.

\subsection{Descomposición en \(5\) modos de polarización}

La dimensi\'on cinco adquiere contenido tensorial al elegir una direcci\'on
unitaria de propagaci\'on \(n\) y un marco ortonormal orientado
\((e_1,e_2,n)\). Definimos los cinco tensores reales
\begin{align}
E_{+}&=\frac1{\sqrt2}(e_1e_1^T-e_2e_2^T),
&E_{\times}&=\frac1{\sqrt2}(e_1e_2^T+e_2e_1^T),
\label{ii:ii:pent:eq:gravity-helicity2-real}\\
E_{1}&=\frac1{\sqrt2}(e_1n^T+ne_1^T),
&E_{2}&=\frac1{\sqrt2}(e_2n^T+ne_2^T),
\label{ii:ii:pent:eq:gravity-helicity1-real}\\
E_{0}&=\frac1{\sqrt6}(2nn^T-e_1e_1^T-e_2e_2^T).
\label{ii:ii:pent:eq:gravity-helicity0-real}
\end{align}
Todos son sim\'etricos y sin traza. Los dos primeros son transversales; los
dos siguientes mezclan el plano transversal con la direcci\'on de
propagaci\'on, y el \`ultimo es el modo escalar longitudinal sin traza.

\begin{theorem}[Descomposici\'on helicoidal completa del portador]
\label{ii:ii:pent:thm:gravity-five-helicities}
Los tensores de
\eqref{ii:ii:pent:eq:gravity-helicity2-real}, \eqref{ii:ii:pent:eq:gravity-helicity1-real}, \eqref{ii:ii:pent:eq:gravity-helicity0-real}
forman una base ortonormal de \(\operatorname{STF}_3\). En la
complejificaci\'on, las combinaciones
\begin{equation}
\label{ii:ii:pent:eq:gravity-complex-helicities}
E_{\pm2}=\frac1{\sqrt2}(E_+\mp i E_\times),
\qquad
E_{\pm1}=\frac1{\sqrt2}(E_1\mp i E_2),
\qquad E_0=E_0
\end{equation}
son vectores de peso \(\pm2,\pm1,0\), respectivamente, bajo rotaciones
alrededor de \(n\). Adem\'as,
\begin{equation}
\label{ii:ii:pent:eq:gravity-tt-on-five}
\Lambda(n)E_{\pm2}=E_{\pm2},
\qquad
\Lambda(n)E_{\pm1}=0,
\qquad
\Lambda(n)E_0=0.
\end{equation}
\end{theorem}

\begin{proof}
En el marco \((e_1,e_2,n)\), cada tensor tiene norma de Frobenius uno.
Sus patrones de entradas no diagonales son disjuntos, y las tres diagonales
de \(E_+,E_0\) tienen producto interno cero; por tanto, los cinco tensores
son ortonormales. Como \(\dim\operatorname{STF}_3=6-1=5\), forman una base.

Una rotaci\'on de \`angulo \(\psi\) alrededor de \(n\) env\'ia
\((e_1,e_2)\) a
\((\cos\psi\,e_1+\sin\psi\,e_2,
-\sin\psi\,e_1+\cos\psi\,e_2)\). Sustituyendo en las f\'ormulas se obtiene
una rotaci\'on de \`angulo \(2\psi\) sobre
\(\operatorname{span}\{E_+,E_\times\}\), una de \`angulo \(\psi\) sobre
\(\operatorname{span}\{E_1,E_2\}\), y \(E_0\) fijo. De ah\'i los pesos de
\eqref{ii:ii:pent:eq:gravity-complex-helicities}. Finalmente,
\(\Pi E_{+}\Pi=E_+\) y \(\Pi E_\times\Pi=E_\times\), mientras
\(\Pi E_1\Pi=\Pi E_2\Pi=0\). Para \(E_0\),
\(\Pi E_0\Pi=-(e_1e_1^T+e_2e_2^T)/\sqrt6\), que es puramente traza en
el plano, y la sustracci\'on de \eqref{ii:ii:pent:eq:gravity-lambda-tt} la anula.
\end{proof}

El teorema precisa el significado de ``cinco polarizaciones'' sin atribuir
a todas el mismo estatuto din\'amico. Antes de imponer ecuaciones de campo,
el portador de esp\'in dos tiene cinco coordenadas irreducibles. En una
realizaci\'on masiva de Fierz--Pauli, las restricciones dejan precisamente
los cinco pesos de \eqref{ii:ii:pent:eq:gravity-complex-helicities}; en una realizaci\'on
sin masa con invariancia de calibre, los sectores \(\pm1\) y \(0\) se eliminan y
quedan \(\pm2\). HMT proporciona aqu\'i el portador, la base y el proyector;
la elecci\'on de la din\'amica masiva o sin masa pertenece a la interfaz
f\'isica. Por eso cinco no significa ``cinco ondas observadas'', sino cinco
grados tensoriales disponibles antes de la reducci\'on.

La orientaci\'on tambi\'en queda visible: cambiar
\((e_1,e_2,n)\) por \((e_1,-e_2,-n)\) conserva \(E_+,E_0\) e intercambia
las fases de las parejas helicoidales. Esa informaci\'on no est\'a en el
escalar \(G_a\); reside en el portador sobre el que act\'ua
\(\mathbb G_{5,a}\).

\subsection{Interpretación física de la cadena representacional}

La dimensi\'on cinco aparece en objetos distintos, y s\'olo
\eqref{ii:ii:pent:eq:gravity-12-5-5-2} autoriza el transporte entre ellos:

\begin{enumerate}
\item \(V_5\) es el sumando irreducible de la representaci\'on sobre los
doce v\'ertices. Es un resultado finito de teor\'ia de representaciones.
\item \(\operatorname{STF}_3\) es el portador real de dimensi\'on cinco de
los pesos \(m=-2,-1,0,1,2\) de esp\'in dos.
\item Una teor\'ia masiva de Fierz--Pauli puede realizar esos cinco pesos
como cinco estados f\'isicos, despu\'es de fijar acci\'on y restricciones.
\item La relatividad general sin masa conserva dos helicidades radiativas,
\(\pm2\), representadas por \(E_+\) y \(E_\times\).
\item Las clases \(E(2)\) de ondas m\'etricas, como \(III_5\), clasifican
posibles amplitudes bajo otras hip\'otesis; no son autom\'aticamente
\(V_5\), Fierz--Pauli o relatividad general.
\end{enumerate}

El contenido estructural queda así determinado: bajo las condiciones del
descenso de incidencia, HMT construye un portador tensorial de cinco
componentes antes de las restricciones. La realización sin masa conserva
las dos helicidades radiativas mediante la proyección TT de rango dos.
La información de orientación se conserva en los tensores y en su
transporte; el valor escalar \(G_a\) determina el acoplamiento común.
